Electrical schematic interpretation requires models to associate terminal labels with conductors, distinguish wiring topology from contact state, and justify conclusions under explicit operating assumptions. Existing visual question answering benchmarks provide limited coverage of this combination in native engineering drawings. This article introduces ELEC-VQA, a source-traceable dataset constructed from computer-aided design drawings of substation secondary circuits. The dataset contains 10,201 English question--answer records across 15 task families and 1,019 image configurations, comprising 300 original local windows and 719 controlled drawing variants. Each record couples a concise answer with compact evidence and executable checks over a bounded terminal graph. Native rendering preserves conductor geometry, while source-disjoint partitions reduce leakage between related questions and edited drawings. Two aligned supervision views enable a controlled comparison of answer-only and evidence-guided fine-tuning of a seven-billion-parameter vision--language model. A comparative study evaluates three general-purpose vision--language model families and two aligned adaptations of a 7B backbone. With provisional reference-based supplementary adjudication of unscored outputs, answer-only fine-tuning increases direct-answer accuracy from 37.16\% to 74.27\%. Under common explanation prompting, evidence-guided fine-tuning achieves 64.24\% joint answer--evidence correctness, compared with 46.17\% for answer-only supervision, while their answer accuracies are 67.97\% and 70.16\%. Evidence support among correct answers and controlled-change analyses reveal selective benefits for multi-condition reasoning and measurement planning, alongside weaknesses in state invariance and visual topology changes. ELEC-VQA enables evaluation of both electrical decisions and their explicit support, without equating generated explanations with faithful internal reasoning.
